% harryopo-book.cls 自检样张（不用于交付，仅验证类文件可编译）
\documentclass[11pt,a4paper]{harryopo-book}

\title{书籍文档类自检样张}
\subtitle{验证五级字体阶梯 / 定理编号 / 教学盒子 / 旁注 / 双面页眉}
\author{harryopo}
\bookseries{自检样张}
\institute{harryopo-office}
\date{2026年9月30日}

\begin{document}
\frontmatter
\maketitle
\tableofcontents

\mainmatter
\chapter{章级结构}
\section{节与五级字体阶梯}
\subsection{小节：黑体}
\subsubsection{小小节：楷体}
\subhead{括号式小标题示例}
正文段落使用方正书宋，正文加粗应使用 {\fzht 黑体} 而非 \textbf。

\bianzhu[旁注]{这是右侧旁注，仿宋小字，用于补充说明。}

\begin{theorem}[勾股定理]
直角三角形两直角边平方和等于斜边平方。
\end{theorem}
\begin{definition}
勾股定理中的三条边分别称为直角边与斜边。
\end{definition}
\begin{remark}
本环境与定义共享计数器。
\end{remark}

公式编号应为章级：
\begin{equation}
  \int_0^1 x^2 \,\mathrm{d}x = \frac{1}{3}
\end{equation}

\chapter{六种教学盒子}
\section{核心概念}
\begin{keyconcept}
向量空间是线性代数的基本对象。
\end{keyconcept}

\section{例}
\begin{examplebox}
求矩阵 $A$ 的特征值。
\end{examplebox}

\section{小贴士}
\begin{tipbox}
计算前先看清题目要求的是行空间还是列空间。
\end{tipbox}

\section{注意}
\begin{warningbox}
相似不等于合同，行列式与迹在相似变换下不变。
\end{warningbox}

\section{练习}
\begin{practicebox}
证明：可逆矩阵的逆矩阵可逆。
\end{practicebox}

\section{解答}
\begin{answerbox}
设 $A^{-1}$ 存在，则 $A^{-1}$ 与 $A$ 的行列式互为倒数，故均非零。
\end{answerbox}

\backmatter
\chapter*{附录}
\end{document}
